% GENERATED FILE. Edit canonical lessons, then run npm run generate:books.
% canonical-source-hash: fnv1a64:ece8fc16096eff5a
% canonical-lessons: ES-C569-luchar, ES-C569-derrotar, ES-C569-vencer, ES-C569-defender, ES-C569-conflicto

\chapter{La lucha}
\label{ch:es-la-lucha}
\hlchaptermodality{569}

\begin{chapteropening}
\textbf{By the end of this chapter:} \emph{I can say that people fought, defended themselves, defeated the enemy and won, and talk about the conflict.}

Complete the last lesson of chapter 569: I can say that people fought, defended themselves, defeated the enemy and won, and talk about the conflict.
\end{chapteropening}

\section[luchar]{luchar — to fight, to struggle}

\label{lesson:ES-C569-luchar}



\begin{quote}
Before the new one: say the Spanish for the queen, then the Spanish for to conquer.
\end{quote}

\subsection*{You'll want to know: luchar}
\textbf{luchar} — \textquotedblleft{}to fight, to struggle\textquotedblright{}.

\begin{quote}
Lucharon por sus derechos.
\end{quote}

\begin{quote}
\emph{They fought for their rights.}
\end{quote}

\emph{Luchar por} — to fight for. \emph{Luchar contra} — to fight against.

\begin{cousinweb}
\emph{luchar} is inherited from Latin \emph{luctārī} 'to wrestle', through spoken Latin \emph{*luctāre} (EtymDB).
\end{cousinweb}

\begin{grammarlens}[title={What you've built}]
One more word for telling, explaining and describing at B1.
\end{grammarlens}

\subsection*{Your turn}
\begin{itemize}
\raggedright
  \item \emph{Say it:} \emph{luchar}
  \item \emph{Say it:} \emph{luchar}, once more
  \item \emph{Say it:} a sentence of your own with \emph{luchar}
  \item \emph{Recall:} say the Spanish for the queen, then the Spanish for to conquer, then say \emph{luchar} again
\end{itemize}

\begin{tcolorbox}[breakable,colback=teal!4,colframe=teal!35!black,title={Before you move on}]
What does luchar mean? (\textquotedblleft{}To fight, to struggle\textquotedblright{}.) Say it once more.
\end{tcolorbox}
\section[derrotar]{derrotar — to defeat}

\label{lesson:ES-C569-derrotar}



\begin{quote}
Before the new one: say the Spanish for to conquer, then the Spanish for to fight.
\end{quote}

\subsection*{You'll want to know: derrotar}
\textbf{derrotar} — \textquotedblleft{}to defeat\textquotedblright{}.

\begin{quote}
Nuestro club derrotó al favorito.
\end{quote}

\begin{quote}
\emph{Our club beat the favourite.}
\end{quote}

\emph{Derrotar a alguien} — to defeat someone.

\begin{cousinweb}
\emph{derrotar} is built on \emph{derrota} (EtymDB), which goes back to Latin \emph{dīrumpere} 'to break apart', from \emph{rumpere} 'to break'.
\end{cousinweb}

\begin{grammarlens}[title={What you've built}]
One more word for telling, explaining and describing at B1.
\end{grammarlens}

\subsection*{Your turn}
\begin{itemize}
\raggedright
  \item \emph{Say it:} \emph{derrotar}
  \item \emph{Say it:} \emph{derrotar}, once more
  \item \emph{Say it:} a sentence of your own with \emph{derrotar}
  \item \emph{Recall:} say the Spanish for to conquer, then the Spanish for to fight, then say \emph{derrotar} again
\end{itemize}

\begin{tcolorbox}[breakable,colback=teal!4,colframe=teal!35!black,title={Before you move on}]
What does derrotar mean? (\textquotedblleft{}To defeat\textquotedblright{}.) Say it once more.
\end{tcolorbox}
\section[vencer]{vencer — to beat, to win; to expire}

\label{lesson:ES-C569-vencer}



\begin{quote}
Before the new one: say the Spanish for to fight, then the Spanish for to defeat.
\end{quote}

\subsection*{You'll want to know: vencer}
\textbf{vencer} — \textquotedblleft{}to beat, to win; to expire\textquotedblright{}.

\begin{quote}
Vencimos por dos a uno.
\end{quote}

\begin{quote}
\emph{We won two-one.}
\end{quote}

It takes \emph{z}: \emph{venzo}. On a document it means to expire: \emph{El plazo vence el lunes.}

\begin{cousinweb}
\emph{vencer} is inherited from Latin \emph{vincere} 'to conquer' (EtymDB), through old Spanish \emph{vencer}.
\end{cousinweb}

\begin{grammarlens}[title={What you've built}]
One more word for telling, explaining and describing at B1.
\end{grammarlens}

\subsection*{Your turn}
\begin{itemize}
\raggedright
  \item \emph{Say it:} \emph{vencer}
  \item \emph{Say it:} \emph{vencer}, once more
  \item \emph{Say it:} a sentence of your own with \emph{vencer}
  \item \emph{Recall:} say the Spanish for to fight, then the Spanish for to defeat, then say \emph{vencer} again
\end{itemize}

\begin{tcolorbox}[breakable,colback=teal!4,colframe=teal!35!black,title={Before you move on}]
What does vencer mean? (\textquotedblleft{}To beat, to win; to expire\textquotedblright{}.) Say it once more.
\end{tcolorbox}
\section[defender]{defender — to defend}

\label{lesson:ES-C569-defender}



\begin{quote}
Before the new one: say the Spanish for to defeat, then the Spanish for to beat.
\end{quote}

\subsection*{You'll want to know: defender}
\textbf{defender} — \textquotedblleft{}to defend\textquotedblright{}.

\begin{quote}
Defendió a su hermano en el colegio.
\end{quote}

\begin{quote}
\emph{He stood up for his brother at school.}
\end{quote}

Stem change: \emph{defiendo}. \emph{Defenderse} — to defend yourself; also to get by.

\begin{cousinweb}
\emph{defender} is inherited from Latin \emph{dēfendere} 'to ward off, to defend' (EtymDB).
\end{cousinweb}

\begin{grammarlens}[title={What you've built}]
One more word for telling, explaining and describing at B1.
\end{grammarlens}

\subsection*{Your turn}
\begin{itemize}
\raggedright
  \item \emph{Say it:} \emph{defender}
  \item \emph{Say it:} \emph{defender}, once more
  \item \emph{Say it:} a sentence of your own with \emph{defender}
  \item \emph{Recall:} say the Spanish for to defeat, then the Spanish for to beat, then say \emph{defender} again
\end{itemize}

\begin{tcolorbox}[breakable,colback=teal!4,colframe=teal!35!black,title={Before you move on}]
What does defender mean? (\textquotedblleft{}To defend\textquotedblright{}.) Say it once more.
\end{tcolorbox}
\section[el conflicto]{el conflicto — the conflict}

\label{lesson:ES-C569-conflicto}



\begin{quote}
Before the new one: say the Spanish for to beat, then the Spanish for to defend.
\end{quote}

\subsection*{You'll want to know: el conflicto}
\textbf{el conflicto} — \textquotedblleft{}the conflict\textquotedblright{}.

\begin{quote}
El conflicto duró muchos años.
\end{quote}

\begin{quote}
\emph{The conflict lasted many years.}
\end{quote}

\begin{cousinweb}
\emph{conflicto} is a learned borrowing from Latin \emph{cōnflīctus} 'a striking together' (Lewis \& Short), from \emph{cōnflīgere} 'to strike together'.
\end{cousinweb}

\begin{grammarlens}[title={What you've built}]
One more word for telling, explaining and describing at B1.
\end{grammarlens}

\subsection*{Your turn}
\begin{itemize}
\raggedright
  \item \emph{Say it:} \emph{el conflicto}
  \item \emph{Say it:} \emph{el conflicto}, once more
  \item \emph{Say it:} a sentence of your own with \emph{el conflicto}
  \item \emph{Recall:} say the Spanish for to beat, then the Spanish for to defend, then say \emph{el conflicto} again
\end{itemize}

\begin{tcolorbox}[breakable,colback=teal!4,colframe=teal!35!black,title={Before you move on}]
What does el conflicto mean? (\textquotedblleft{}The conflict\textquotedblright{}.) Say it once more.
\end{tcolorbox}
